% STATUS: written 2026-10-21, recalled results as is (one environment each, source kind and whole statement)
% STATUS: extended 2026-11-06 (Ghosh and Prakash 2018 and Vladeva 2018 read in full from the library PDFs)
% Sources read in full: Cheng and Jing 2007 (arXiv, pp. 1--5), Bresar 2005 (pp. 411--425, Sections 4 and 5); Ghosh and Prakash 2018 (pp. 1--5, 8); Vladeva 2018 (pp. 1--5).
\section{Jordan derivations of specific algebras}\label{sec:jordan-type:jordan-on-algebras}

Section~\ref{sec:jordan-type:semiprime} gave conditions on a ring that force a Jordan derivation to be
a derivation. Cheng and Jing study triangular algebras \cite[p.~1]{ChengJing2007}. Brešar surveys a construction of
Benkovič for upper triangular matrices \cite[\S4.1, p.~418]{Bresar2005}. He also treats $C^{*}$-algebras
\cite[Section~5, p.~421]{Bresar2005}.

Let $A$ and $B$ be algebras over a commutative ring $R$ \cite[p.~2]{ChengJing2007}. Let $M$ be an
$(A,B)$-bimodule that is faithful as a left $A$-module and as a right $B$-module
\cite[p.~2]{ChengJing2007}. These data give the triangular algebra $T=\mathrm{Tri}(A,M,B)$, whose elements are the
matrices $\begin{pmatrix}a&m\\0&b\end{pmatrix}$ \cite[p.~2]{ChengJing2007}.
Zhang and Yu showed that every Jordan derivation of a triangular algebra is a derivation
\cite[p.~1]{ChengJing2007}. Cheng and Jing state it as follows \cite[p.~1]{ChengJing2007}.

\begin{theorem}[{\cite[Theorem~1.1, p.~1]{ChengJing2007}}]\label{thm:jordan:zhang-yu}
Let $A,B$ be unital algebras over a $2$-torsion free commutative ring $R$, and $M$ be a unital
$(A,B)$-bimodule that is faithful as a left $A$-module and also as a right $B$-module. Then every
Jordan derivation from the triangular algebra $\mathrm{Tri}(A,M,B)$ into itself is a derivation.
\end{theorem}

The proof of Zhang and Yu needs both algebras to have a unit \cite[p.~1]{ChengJing2007}.
Cheng and Jing do not require a unit in either algebra \cite[p.~1]{ChengJing2007}. They assume the following property of $A$ and $B$
\cite[Section~2, p.~2]{ChengJing2007}: (P) if $a\in A$ (respectively $B$) and $xay+yax=0$ for all
$x,y\in A$ (respectively $B$), then $a=0$ (Theorem~\ref{thm:jordan:cheng-jing}).

\begin{theorem}[{\cite[Theorem~2.10, p.~5]{ChengJing2007}}]\label{thm:jordan:cheng-jing}
Let $A$ and $B$ be two algebras over a $2$-torsion free commutative ring $R$ with the property:
(P) Suppose that $a\in A$ (respectively $B$). If $xay+yax=0$ holds for all $x,y\in A$ (respectively $B$),
then $a=0$. Let $M$ be a faithful $(A,B)$-bimodule and $T$ be the triangular algebra
$\mathrm{Tri}(A,M,B)$. Then every Jordan derivation $\delta$ on $T$ into itself is a derivation.
\end{theorem}

Cheng and Jing note that Theorem~\ref{thm:jordan:zhang-yu} is the special case where both
algebras are unital \cite[Remark~2.11, p.~5]{ChengJing2007}. Their proof is a sequence of lemmas.
The first shows that $\delta$ is a derivation on each of the three parts $T_{11}$, $T_{12}$ and $T_{22}$
of $T$ \cite[pp.~4--5]{ChengJing2007}.

Now let $F$ be a field with
$\operatorname{char}(F)\neq2$, and let $T_n(F)$ be the algebra of upper triangular $n\times n$
matrices over $F$ \cite[Section~4, p.~418]{Bresar2005}. Brešar describes a construction of
Benkovič \cite[\S4.1, p.~418]{Bresar2005}. It produces maps $d$ with $d(xy)=d(y)x+yd(x)$ for all $x,y$, which are called
antiderivations \cite[\S4.1, p.~418]{Bresar2005}. The construction gives a proper Jordan derivation
on $T_n(F)$ into a bimodule \cite[\S4.1, p.~418]{Bresar2005}.

\begin{example}[{\cite[Example~4.1, p.~418]{Bresar2005}}]\label{ex:jordan:benkovic-tn}
Let $A=T_n(F)$, $n\geqslant2$. Then there exists an $A$-bimodule $M$ and a proper Jordan derivation
(in fact, an antiderivation) $d\colon A\to M$ such that $d(ue)\neq d(u)e+ud(e)$ for some
$u\in A^{(1)}$ and some idempotent $e\in A$.
\end{example}

Brešar adds that the main result of Benkovič says that every Jordan derivation from $T_n(F)$ into a
$T_n(F)$-bimodule is the sum of a derivation and an antiderivation \cite[p.~419]{Bresar2005}. So
Jordan derivations of $T_n(F)$ act as derivations on a large part of the algebra
\cite[p.~419]{Bresar2005}.

For operator algebras, Brešar poses a question: is every Jordan derivation from a
$C^{*}$-algebra $A$ into a Banach $A$-bimodule a derivation \cite[Section~5, p.~421]{Bresar2005}?
He calls it an intriguing open question \cite[Section~5, p.~421]{Bresar2005}. Johnson proved that a
continuous Jordan derivation from a $C^{*}$-algebra into a Banach bimodule is a derivation
\cite[p.~421]{Bresar2005}. For additive maps Brešar gives a definitive answer on the algebras he
calls $\mathrm{CM}$-algebras \cite[p.~421]{Bresar2005}.

\begin{theorem}[{\cite[Theorem~5.1, p.~422]{Bresar2005}}]\label{thm:jordan:bresar-cm}
Let $A$ be a unital $\mathrm{CM}$-algebra. Then there exists a proper additive Jordan derivation
from $A$ into some unital $A$-bimodule if and only if $A$ contains an ideal of codimension $1$.
\end{theorem}

Brešar derives the theorem of Johnson from Theorem~\ref{thm:jordan:bresar-cm}
\cite[p.~422]{Bresar2005}. He also notes that the result of Johnson does not extend from
$\mathrm{CM}$-algebras to semisimple Banach algebras, because a primitive Banach algebra can have
$T_2(\mathbb{C})$ as a quotient \cite[p.~422]{Bresar2005}.

The $C^{*}$-algebra question of this section is open, and this section states
it as a problem. % claim-ok: own signpost introducing the moved problem, which carries its own citation

\begin{problem}[{\cite[Section~5, p.~421]{Bresar2005}}]\label{prob:open:jordan-cstar}
Let $A$ be a $C^{*}$-algebra and let $M$ be a Banach $A$-bimodule.
Is every Jordan derivation from $A$ into $M$ a derivation?
\end{problem}

Bre\v{s}ar calls this an intriguing open question
\cite[Section~5, p.~421]{Bresar2005}. Johnson had proved it for continuous
Jordan derivations, and Bre\v{s}ar answers it for the algebras he calls
$\mathrm{CM}$-algebras \cite[p.~421]{Bresar2005}. The statement of
Theorem~\ref{thm:jordan:bresar-cm} is the answer on that class.
Over a commutative ring with unity the torsion hypothesis can be dropped
\cite[Theorem~2.1, p.~3]{GhoshPrakash2018}. Ghosh and Prakash prove that every
Jordan derivation of the upper triangular matrix algebra $T_n(C)$ into itself is
a derivation, with no condition on the ring
\cite[Theorem~2.1, p.~3]{GhoshPrakash2018}. Their proof computes the image
of each matrix unit and compares coefficients
\cite[pp.~3--4]{GhoshPrakash2018}. The same method shows that every derivation of
$T_n(C)$ is inner \cite[Theorem~2.2, p.~5]{GhoshPrakash2018}. For the 2-element
field the result covers a ring that is not 2-torsion free
\cite[Corollary~2.1, p.~4]{GhoshPrakash2018}. The full matrix algebra behaves in
the same way \cite[Theorem~3.1, p.~5]{GhoshPrakash2018}. Every Jordan derivation
of $M_n(C)$ into itself is a derivation, and every derivation of $M_n(C)$ is
inner \cite[Theorem~3.2, p.~8]{GhoshPrakash2018}.

The passage from rings to semirings keeps the triangular shape, and it weakens
the conclusion. % claim-ok: transition to the semiring case, the plan of the book
Every derivation of the upper triangular matrix semiring
keeps all entries of a matrix except those of one principal dense submatrix,
and it sets those entries to zero
\cite[Theorem~1, p.~5]{Vladeva2018Semirings}. Vladeva works over an additively
idempotent semiring, where the sum of two equal terms adds nothing
\cite[Proposition~1, p.~3]{Vladeva2018Semirings}. In that setting the Jordan
product of a matrix with a partial identity matrix equals the ordinary product,
so the map $A\mapsto A\circ D_k$ is a derivation
\cite[Proposition~1, p.~3]{Vladeva2018Semirings}. The paper counts the
derivations that arise this way \cite[Theorem~1, p.~5]{Vladeva2018Semirings}.

On matrix rings the Jordan condition can be tested through products alone.
Ghahramani characterizes the Jordan derivations of a matrix ring by their action
on zero products \cite[Abstract]{Ghahramani2015ZeroJordan}.

For prime rings the generalized notion has its own theory. Jing and Lu initiate
the study of generalized Jordan derivations and generalized Jordan triple
derivations on prime rings and on standard operator algebras
\cite[Abstract, p.~605]{JingLu2003GeneralizedJordan}.

